\section{Zadanie 7}
\autor{Aleksander Wieczorek}
W samolocie leci \(N = n_1+n_2+n_3+n_4+1\) rodzin, gdzie \(n_i\) rodzin ma \(i\) sztuk
bagażu, dla \(i \in \{1,2,3,4\}\), a nasza rodzina jest jedyną z pięcioma bagażami.\\
Przyjmując, że kolejność wydawania walizek jest losowa, jaka jest oczekiwana liczba rodzin, które będą
czekać dłużej od naszej na swoje bagaże?

\subsection*{Rozwiązanie}
Jak zwykle rozbijemy problem na indykatory. Oznaczmy:\\
\[X= \text{Liczba rodzin czekająca dłużej od nas}\]
\[X_{ij}=\begin{cases}
		1 & \text{jeśli j-ta rodzina z \(n_i\) czeka dłużej od nas} \\
		0 & \text{w p.p.}
	\end{cases}
\]
Żeby policzyć \(P(X_{ij})\) będziemy rozważać tylko jak są ustawione bagaże naszej rodziny i rodziny, która ma i bagaży (względem siebie).\\
Tzn. rozpatrujemy tylko sytuacje, gdy mamy 5+i bagaży oraz 1 spośród i bagaży wyjedzie ostatni.\\
\includegraphics[width=\textwidth]{./tasks/07.png}
W powyższym przykładzie niebieskie kwadraty to nasze walizki, a czerwone to bagaże jakiejś rodziny z \(n_3\)\\
Ostatni bagaż jest ustalony, więc ustawień dla których \(X_{3j}=1\) jest \(\binom{7}{2}\), a wszystkich jest \(\binom{8}{3}\), stąd \(P(X_{3j})=\frac{\binom{7}{2}}{\binom{8}{3}}\)\\
Ogólniej dla każdej rodziny z \(n_i\) mamy \(P(X_{ij})=\frac{\binom{5+i-1}{i-1}}{\binom{5+i}{i}}\)\\
Liczymy wartość oczekiwaną jak zawsze:\\
\[E(X)=E(\sum_i \sum_j X_{ij})=\sum_i \sum_jE(X_{ij})=\sum_{i=1}^4 n_i\cdot \frac{\binom{5+i-1}{i-1}}{\binom{5+i}{i}}\]
Teraz możemy uprościć \(\frac{\binom{5+i-1}{i-1}}{\binom{5+i}{i}}\) do \(\frac{i}{i+5}\).\\
Zapisujemy mianownik jako \(\frac{i+5}{i} \binom{5+i-1}{i-1}\), dwumian się skróci, więc dostaniemy:
\subsection*{Odpowiedź}
\[E(X)=\sum_{i=1}^4 n_i \cdot \frac{i}{i+5}\]
